\documentclass[12pt]{article}
\usepackage{graphicx}
\usepackage{hyperref}
\usepackage[a4paper, total={6in, 8in}]{geometry}
\usepackage{setspace}
\usepackage{parskip}
\usepackage{tabularx}
\usepackage{multirow}
\usepackage{booktabs}
\usepackage{amsmath}
\usepackage{listings}
\usepackage{xcolor}
\usepackage{float}

\linespread{1.1}

\lstset{
  basicstyle=\ttfamily\small,
  keywordstyle=\color{blue},
  commentstyle=\color{gray},
  stringstyle=\color{red},
  breaklines=true,
  frame=single,
  numbers=left,
  numberstyle=\tiny\color{gray}
}

\title{Training GPT-2 on SEC-EDGAR Financial Filings:\\
A 124M Parameter Language Model for Corporate Disclosure Generation}
\author{Li Zhiwei}
\date{June 2026}

\begin{document}

\maketitle

\begin{abstract}
We present a 124-million parameter GPT-2 language model trained on 1.55 billion tokens of SEC-EDGAR financial filings using the nanoGPT framework. The model was trained on a single NVIDIA RTX 4070 GPU over approximately 5 days, reaching a validation loss of 2.45 at step 37,000 of 47,400 total steps (78\% completion). We evaluate the model's generation quality across multiple SEC filing sections including business descriptions, management discussion and analysis, risk factors, financial notes, and proxy statements. Our analysis reveals that the model successfully learns SEC document structure, financial vocabulary, and boilerplate language patterns, but exhibits characteristic limitations in long-range coherence, numerical consistency, and table extension. We discuss the implications for using small language models in financial document generation and the trade-offs between model size, training data quality, and generation fidelity.
\end{abstract}

\section{Introduction}

The proliferation of large language models (LLMs) has created significant interest in domain-specific text generation. Financial filings submitted to the U.S. Securities and Exchange Commission (SEC) through the EDGAR system represent a particularly rich corpus of structured, domain-specific text. These documents follow strict formatting conventions and employ specialized financial vocabulary, making them an interesting testbed for evaluating language model capabilities in structured document generation.

While state-of-the-art models like GPT-4 and Claude demonstrate impressive general-purpose capabilities, the cost and computational requirements of training and deploying such models limit their accessibility. In this work, we investigate whether a smaller, more accessible model---a 124M parameter GPT-2---can learn to generate convincing SEC filing text when trained on a focused corpus of EDGAR documents.

Our contributions include:
\begin{itemize}
    \item Training a 124M parameter GPT-2 model on 1.55 billion tokens of SEC-EDGAR filings using consumer hardware
    \item Systematic evaluation of generation quality across five distinct SEC filing sections
    \item Analysis of long-range coherence, numerical consistency, and structural fidelity
    \item Discussion of practical applications and limitations for financial document generation
\end{itemize}

\section{Related Work}

\subsection{Language Models for Financial Text}

Recent work has explored applying language models to financial domains. BloombergGPT \cite{bloomberggpt} trained a 50B parameter model on a mixed corpus including financial data. FinGPT \cite{fingpt} proposed open-source financial language models. However, these efforts focused on much larger models and broader financial tasks.

\subsection{Small Language Models}

The original GPT-2 \cite{gpt2} demonstrated that language models could generate coherent text at various scales (124M to 1.5B parameters). More recent work has shown that smaller models can achieve surprising performance when trained on high-quality, domain-specific data \cite{tinystories}.

\subsection{SEC-EDGAR as Training Data}

The SEC-EDGAR database contains millions of corporate filings spanning decades. Prior work has used EDGAR data for sentiment analysis \cite{edgar-sentiment} and information extraction \cite{edgar-extraction}, but little work has explored using EDGAR as a primary training corpus for generative language models.

\section{Methodology}

\subsection{Data Collection and Preparation}

We collected SEC-EDGAR filings from the CodeParrot GitHub dataset, which contains processed EDGAR documents. The training corpus consists of:

\begin{table}[H]
\centering
\begin{tabular}{lr}
\toprule
\textbf{Metric} & \textbf{Value} \\
\midrule
Total training tokens & 1.55B \\
Number of training shards & 16 \\
Validation tokens & 100M \\
Average document length & ~2,048 tokens \\
Vocabulary & GPT-2 BPE (50,257 tokens) \\
\bottomrule
\end{tabular}
\caption{SEC-EDGAR training corpus statistics}
\label{tab:corpus}
\end{table}

The data was tokenized using the GPT-2 byte-pair encoding tokenizer and stored as memory-mapped numpy arrays for efficient training.

\subsection{Model Architecture}

We use the standard GPT-2 architecture with the following configuration:

\begin{table}[H]
\centering
\begin{tabular}{lr}
\toprule
\textbf{Parameter} & \textbf{Value} \\
\midrule
Number of layers & 12 \\
Number of attention heads & 12 \\
Embedding dimension & 768 \\
Context length & 1,024 tokens \\
Total parameters & 123.59M \\
Vocabulary size & 50,257 \\
\bottomrule
\end{tabular}
\caption{GPT-2 124M model configuration}
\label{tab:model}
\end{table}

\subsection{Training Setup}

Training was conducted on a single NVIDIA RTX 4070 GPU with 12GB VRAM using the nanoGPT framework:

\begin{table}[H]
\centering
\begin{tabular}{lr}
\toprule
\textbf{Training Parameter} & \textbf{Value} \\
\midrule
Batch size (micro) & 4 \\
Gradient accumulation steps & 8 \\
Effective batch size & 32,768 tokens \\
Learning rate & 6e-4 \\
Minimum learning rate & 6e-5 \\
Warmup iterations & 2,000 \\
Maximum iterations & 47,400 \\
Weight decay & 0.1 \\
Optimizer & AdamW ($\beta_1=0.9$, $\beta_2=0.95$) \\
Gradient clipping & 1.0 \\
\bottomrule
\end{tabular}
\caption{Training hyperparameters}
\label{tab:training}
\end{table}

Training required approximately 5 days to reach step 37,000 (78\% completion), achieving a validation loss of 2.45.

\section{Results}

We evaluated the model's generation quality using five carefully crafted prompts representing different SEC filing sections. Each prompt contained 500--1000 characters of authentic SEC text, and we generated 1,200 tokens of continuation at temperature 0.7.

\subsection{Prompt 1: Business Description (Item 1)}

\textbf{Input:} A healthcare SaaS company description with \$487.2M revenue, 8,200 employees, and EHR platform details.

\textbf{Observations:}
\begin{itemize}
    \item First 2--3 paragraphs maintained healthcare industry context
    \item Gradually shifted from SaaS to biotech/pharma framing
    \item After ~200 generated tokens, entered a ``commercialization of product candidates'' loop
    \item Hallucinated drug names: X-Avent, X-Zentib, S-Zentib, Q-partnerib
    \item Maintained grammatical structure throughout, even in degenerate sections
\end{itemize}

\subsection{Prompt 2: MD\&A (Revenue Analysis)}

\textbf{Input:} Revenue growth of 28\%, cost analysis, and gross margin improvement from 62.7\% to 64.4\%.

\textbf{Observations:}
\begin{itemize}
    \item First continuation paragraph was excellent: ``Cost of revenue increased by \$32.1 million, or 26\%... primarily attributable to decreased depreciation and amortization expense in the period of the acquisition of DMR''
    \item Subsequently generated 10 consecutive paragraphs all starting with ``Cost of revenue increased/decreased by \$X.X million''
    \item Numbers became internally inconsistent (e.g., ``decreased by \$61.2 million, or 13\%, to \$1.9 million from \$2.7 million'')
    \item Percentage calculations frequently contradicted the dollar amounts cited
\end{itemize}

\subsection{Prompt 3: Risk Factors}

\textbf{Input:} Net losses of \$42.3M, \$67.8M, and \$89.1M over three years, with \$523.4M accumulated deficit.

\textbf{Observations:}
\begin{itemize}
    \item Maintained risk factor structure (heading + explanation) for extended generation
    \item ``Product candidates'' appeared 47 times in 25 lines
    \item Recursive self-reference: ``Our product candidates may fail to develop, develop and commercialize our product candidates may fail''
    \item Lost thread of healthcare SaaS, defaulted to generic biotech risk factors
\end{itemize}

\subsection{Prompt 4: Financial Notes (Revenue Recognition)}

\textbf{Input:} Revenue disaggregation table with subscription (\$380M), professional services (\$89M), and hardware (\$18M) revenue.

\textbf{Observations:}
\begin{itemize}
    \item Table echoed with perfect formatting and alignment
    \item All numbers reproduced exactly as in the input
    \item Attempted to start a second table but generated only blank spaces
    \item Only ~50 tokens of meaningful continuation before collapse
\end{itemize}

\subsection{Prompt 5: Proxy Statement (Executive Compensation)}

\textbf{Input:} Executive compensation table with three named officers and full compensation breakdown.

\textbf{Observations:}
\begin{itemize}
    \item Table structure perfectly preserved with correct column alignment
    \item All three executive rows echoed exactly
    \item Attempted to add a fourth executive (``William R. Gras'') but only generated partial row
    \item Immediately collapsed to blank spaces after incomplete row
\end{itemize}

\section{Analysis}

\subsection{Echo vs. Generate Distinction}

A clear pattern emerges: the model excels at \textbf{echoing} input content (tables, numbers, formatting) but struggles to \textbf{generate} new content that maintains consistency. This suggests the model learned surface-level patterns rather than underlying data relationships.

\subsection{Loop Attractors}

Certain phrases act as ``probability sinks'' that the model falls into when uncertain:
\begin{itemize}
    \item ``Commercialization of product candidates'' (Prompts 1, 3)
    \item ``Cost of revenue increased by \$X million'' (Prompt 2)
    \item ``Raise additional capital'' (observed in earlier tests)
\end{itemize}

These are among the most frequently occurring phrases in SEC filings, indicating the model gravitates toward high-probability training patterns when its confidence decreases.

\subsection{Numerical Coherence}

\begin{itemize}
    \item \textbf{Dollar amounts:} Plausible scale (\$1M--\$500M range) but internally inconsistent
    \item \textbf{Percentages:} Often don't match the dollar changes cited
    \item \textbf{Dates:} Consistent (always ``December 31, 2023/2022'')
    \item \textbf{Implication:} The model learned the \textit{format} of numbers, not their \textit{meaning}
\end{itemize}

\subsection{Domain Drift}

The healthcare SaaS prompt drifted to biotech/pharma within 200 tokens, suggesting:
\begin{itemize}
    \item Training data may be dominated by biotech 10-K filings
    \item Biotech risk factors serve as ``default'' SEC content when the model is uncertain
    \item Domain-specific fine-tuning data distribution significantly affects generation quality
\end{itemize}

\subsection{Grammar vs. Logic}

Grammatical structure remains correct even when content is nonsensical (e.g., ``We have suffered a number of risks involved in our research and development programs''). Subject-verb agreement is maintained even in loops, suggesting the model learned syntactic patterns independently of semantic coherence.

\section{Discussion}

\subsection{Practical Applications}

The 124M SEC-EDGAR model demonstrates capability for:
\begin{itemize}
    \item Generating SEC boilerplate language and standard disclosures
    \item Suggesting section structures and formatting conventions
    \item Drafting placeholder text that visually resembles authentic filings
    \item Providing a starting point for human writers to edit and refine
\end{itemize}

\subsection{Limitations}

The model is \textbf{not} suitable for:
\begin{itemize}
    \item Generating accurate financial data or calculations
    \item Maintaining consistency across long documents
    \item Producing factually grounded content
    \item Extending tables with new, coherent rows
\end{itemize}

\subsection{Comparison with Larger Models}

These limitations are consistent with observations of GPT-2 1.5B (12$\times$ larger), which shows similar but less severe patterns. The 124M model's behavior aligns with expectations for its parameter count and suggests that:
\begin{itemize}
    \item Document structure and vocabulary can be learned at small scales
    \item Numerical reasoning requires significantly more parameters or specialized architectures
    \item Long-range coherence improves with scale but remains challenging
\end{itemize}

\subsection{Future Work}

Potential improvements include:
\begin{enumerate}
    \item \textbf{Data quality:} Curating a balanced corpus across filing types and industries
    \item \textbf{Model architecture:} Incorporating numerical embeddings or structured prediction heads
    \item \textbf{Training objectives:} Adding auxiliary losses for numerical consistency
    \item \textbf{Post-processing:} Implementing verification layers for generated numbers
    \item \textbf{Scale:} Training larger models (350M, 1.3B) on the same corpus
\end{enumerate}

\section{Conclusion}

We have demonstrated that a 124M parameter GPT-2 model can learn SEC-EDGAR document structure, financial vocabulary, and boilerplate language when trained on 1.55 billion tokens of corporate filings. The model achieves convincing surface-level generation quality but exhibits fundamental limitations in numerical consistency, long-range coherence, and table extension.

These results suggest that small language models can serve as useful tools for financial document drafting assistance, particularly for generating boilerplate sections and suggesting document structures. However, they cannot replace human expertise for ensuring accuracy and consistency in actual SEC filings.

The SEC-EDGAR corpus presents unique challenges for language models due to its structured nature, numerical density, and strict formatting requirements. Our work provides a baseline for future research on domain-specific financial language models and highlights the trade-offs between model accessibility and generation fidelity.

\begin{thebibliography}{9}

\bibitem{bloomberggpt}
Wu, S., et al. (2023).
\textit{BloombergGPT: A Large Language Model for Finance}.
arXiv preprint arXiv:2303.17564.

\bibitem{fingpt}
Yang, H., et al. (2023).
\textit{FinGPT: Open-Source Financial Large Language Models}.
arXiv preprint arXiv:2306.06031.

\bibitem{gpt2}
Radford, A., et al. (2019).
\textit{Language Models are Unsupervised Multitask Learners}.
OpenAI Technical Report.

\bibitem{tinystories}
Eldan, R., \& Li, Y. (2023).
\textit{TinyStories: How Small Can Language Models Be and Still Speak Coherent English?}
arXiv preprint arXiv:2305.07759.

\bibitem{edgar-sentiment}
Loughran, T., \& McDonald, B. (2011).
\textit{When Is a Liability Not a Liability? Textual Analysis, Dictionaries, and 10-Ks}.
The Journal of Finance, 66(1), 35-65.

\bibitem{edgar-extraction}
Li, Y., et al. (2020).
\textit{Deep Learning for Financial Document Information Extraction}.
Proceedings of the 1st Joint Workshop on Financial Narrative Processing and MultiLing Financial Summarisation.

\end{thebibliography}

\end{document}
